% Generated draft: requires mathematical and source review.
\documentclass[11pt]{article}
\usepackage[margin=25mm]{geometry}
\usepackage{fontspec}
\setmainfont{Latin Modern Roman}
\usepackage{amsmath,amssymb,mathtools,unicode-math}
\setmathfont{Latin Modern Math}
\usepackage{enumitem,needspace,xurl,hyperref}
\setlength{\parindent}{0pt}
\setlength{\parskip}{5pt}
\setlength{\emergencystretch}{3em}
\newcommand{\N}{\mathbb{N}}
\newcommand{\Z}{\mathbb{Z}}
\newcommand{\Q}{\mathbb{Q}}
\newcommand{\R}{\mathbb{R}}
\newcommand{\C}{\mathbb{C}}
\newcommand{\F}{\mathbb{F}}
\newcommand{\A}{\mathbb{A}}
\newcommand{\PP}{\mathbb{P}}
\newcommand{\E}{\mathbb{E}}
\newcommand{\eps}{\varepsilon}
\newcommand{\dd}{\,\mathrm{d}}
\begin{document}
\section*{UnsolvedMath ID: 20002315}
\textbf{Source problem code:} catalogue entry
\subsection{Definitions and mathematical statement}
For a complete theory $T$ in a countable language, write $I(T,\kappa)$ for the number of isomorphism types of models of cardinality $\kappa$.

\textbf{Question 13 and absent antecedent.} ``What if $T$ has finitely many models in $\aleph_1$? (Maybe look at models in $\aleph_n$.)'' The source places this at the end of Section 1.6, ``$\lambda$-many models of each cardinality $\lambda\ge\aleph_1$ (Greenberg)''. The preceding property being modified is not included. The aid interprets the question in the surrounding computable-presentation setting, asking what finite uncountable spectra imply for presentations of the individual models; this interpretation and its effectiveness assumptions need the missing passage, not silent invention.

For a precise convention on the supplied examples, a $\kappa$-computable presentation has domain $\kappa$ and atomic diagram $\Delta_1$ over $L_\kappa$, in the stated computable language and with permitted parameters specified. Relative or parameter-free variants are different; an arbitrary non-effective theory cannot automatically be inserted into an effective presentation theorem.

\textbf{Supplied conditional localization.} The aid invokes Lachlan's theorem and the Greenberg--Knight effective saturated-model construction to say that, for complete countable $T$ with finite $I(T,\aleph_1)$, a possible counterexample in the intended framework is confined to a non-saturated $\aleph_1$-sized model of an $\aleph_0$-categorical but non-$\aleph_1$-categorical theory. This is a reduction, not an answer for that remaining branch. Full hypotheses on the theory, computability base and the construction are not supplied, so it is not asserted here as an unconditional effective result for arbitrary countable theories. Saturated means realizing every type over a parameter set of size less than the model's cardinality.

\textbf{Supplied exact finite-spectrum family.} For $r\ge1$, let $T_r$ be the theory of a partition by $r$ labeled unary predicates $U_1,\ldots,U_r$, each infinite, with no further structure. For each finite $n\ge1$, its $\aleph_n$-sized model types are indexed by
\[
 (|U_1|,\ldots,|U_r|)
  =(\aleph_{i_1},\ldots,\aleph_{i_r}),\qquad
 i_j\in\{0,\ldots,n\},\quad\max_j i_j=n.
\]
Since the labels are named, no permutation quotient is taken. Hence
\[
 I(T_r,\aleph_n)=(n+1)^r-n^r,
 \qquad I(T_r,\aleph_1)=2^r-1.
\]
The aid asserts that every such type has a uniform $\aleph_n$-computable presentation by bounded definitions over $L_{\aleph_n}$ from the finite cardinal-label data, in its allowed-parameter convention. This demonstrates a concrete benign finite-spectrum family; it does not classify all finite-spectrum theories or prove the desired presentation property in the remaining non-saturated branch.

\textbf{Separate Section 1.7, ``Non-abelian free groups (Knight)''.} The appended paragraph concerns a finite family $H_{n,t}$ of $n$-generated one-relator groups with relator length at most $t$, and defines, for a sentence $\varphi$,
\[
 h_{n,t,\varphi}=
   \frac{|\{G\in H_{n,t}:G\models\varphi\}|}{|H_{n,t}|}.
\]
The sampling convention, such as presentations versus marked groups versus isomorphism classes, is not specified. No limit or decidability question for these frequencies occurs before the extract ends. This unrelated group paragraph is not a hypothesis of Question 13.
\subsection{Short English statement}
The source's finite-spectrum question is missing the property it asks to extend. The aid supplies a conditional reduction for computable presentations and an exact labeled-unary-part family with $(n+1)^r-n^r$ models in $\aleph_n$, all explicitly presentable in its convention. The unrelated one-relator frequency definition is retained as separate context.
\subsection{Sources}
\begin{itemize}
\item Local UnsolvedMath registry, permanent ID 20002315, source code catalogue entry; canonical \texttt{aim-logic-notes.json}, notes[90].
\item AIM \emph{Computable stability theory}, Question 13, Section 1.6 (Greenberg), followed by Section 1.7 (Knight): \url{https://aimath.org/pastworkshops/computestabproblems.pdf}.
\item Lachlan's theorem and Greenberg--Knight effective saturated-model construction are named in the aid without full citations, precise hypotheses or computability-parameter conventions. The record reports a PDF boundary check, not performed here.
\item Registry research attempt 1, dated 2026-08-14: partial result, status at run unknown, rigor ``proved'', novelty confidence low. The explicitly machine-generated, non-peer-reviewed aid requires independent verification before citation.
\end{itemize}
\end{document}
